\documentclass[12pt]{article}

\usepackage[utf8]{inputenc}
\usepackage[T1]{fontenc}

\title{La linguistica de la biologia}
\author{Dana Barandica,Junior Sabogal, Dulce Velasquez y Miguel Lopez}
\date{\today}

\begin{document}

\maketitle

\section{Introducción}
<<<<<<< HEAD



A partir de esta relacion surge la biolinguistica, la disciplina que estudia las bases entre la biologia y la linguistica, su objetivo es comprender como se desarrolla la capacidad, que estructuras participan, que influencia tienen en los genes y como aparecio el lenguaje durante la evolución de nuestra especie (Boeckx & Martins, 2016)



\section{Desarrollo}

Ademas del cerebro, la genetica tiene un papel fundamental en el desarrollo de las habilidades linguisticas. Los genes que participan en la formación y funcionamiento de las conexiones de



\section{Conclusiones}

\section{Referencias bibliográficas}

\end{document}
=======
El lenguaje es una de las capacidades mas importantes del ser humano, ya quue permite expresar pensamientos, compartir conociminetos y comunicarse con otras personas. Aunque generalmente se estudia desde el punto de vista social y cultural, tambien tiene una relacon muy importante con la biologia. La capacidad de hablar; comprender y aprender un idioma depende de diferentes procesos que ocurren en nuestro organismo, especialmente en el cerebro y el sistema nervioso.



\end{document}
>>>>>>> b6eba9ebb5fc7e865e968d097e2d16ab60c68144
